\subsection{Summary of Theory Developments} 
In this section, we shortly collect results available for the theoretical description of the associated production of a Higgs boson and a weak vector boson (VH) at the LHC, and we outline recent developments in theory predictions. We recall that the dominant partonic contribution to the process comes from the $q\bar{q}$-initiated, Drell-Yan-like channel, while a subleading contribution due to a loop-induced $gg$-initiated channel arises at NNLO QCD for the $ZH$ final state.

Concerning fixed-order predictions, the QCD corrections for the $q \bar{q}$-initiated process have been known at NNLO~\cite{HAN1991167, Brein:2003wg}. For the Drell-Yan-like contribution, N$^3$LO QCD corrections have become available~\cite{Baglio:2022wzu}. Electroweak (EW) NLO corrections have been computed in Refs.~\cite{Ciccolini:2003jy, Denner:2011id} and are available in the \texttt{HAWK} code~\cite{Denner:2011id,Denner:2014cla}, including the leptonic decays of the weak bosons. Differential predictions are available in Refs.~\cite{Ferrera:2011bk, Ferrera:2013yga,Ferrera:2014lca,Campbell:2016jau,Gauld:2021ule}, including NNLO QCD and NLO EW effects.
In recent years, considerable interest has been devoted to the QCD corrections to the subdominant process $gg \to ZH$~\cite{Dicus:1988yh,Kniehl:1990iva}, with the aim of reducing the non-negligible impact of the associated scale uncertainties. The NLO QCD corrections have been obtained using multiple approaches~\cite{Altenkamp:2012sx,Hasselhuhn:2016rqt,Davies:2020drs,Chen:2020gae,Wang:2021rxu,Alasfar:2021ppe,Chen:2022rua,Degrassi:2022mro,Davies:2025out}. The Refs.~\cite{Chen:2022rua,Degrassi:2022mro} discuss an additional source of theoretical uncertainty due to the choice of renormalization scheme for the top-quark mass ($m_t$). More recently, QCD corrections at NNLO have been computed~\cite{Davies:2025otz} in the large-$m_t$ limit. 

Given the relevance of VH production as an experimental probe of the Higgs coupling to bottom quarks, several works have discussed the effects of supplementing VH predictions with the Higgs decay to a $b \bar{b}$ pair~\cite{Caola:2017xuq,Gauld:2019yng,Behring:2020uzq}.

The resummation of threshold logarithms for the Drell-Yan (DY) component of VH production has been carried out at N$^3$LL accuracy~\cite{Kumar:2014uwa, Das:2022zie}. Soft-gluon resummation in the $gg$-initiated channel has been studied in Refs.~\cite{Harlander:2014wda,Das:2025wbj,Das:2025qym}. The matching of fixed-order corrections to the parton shower (PS) has been obtained for VH production at NLO QCD~\cite{Luisoni:2013cuh,Granata:2017iod} and at NLO EW~\cite{Granata:2017iod} accuracy. Considering NNLO QCD corrections, merged predictions including one jet for $gg \to ZH$ have been presented~\cite{Hespel:2015zea, Goncalves:2015mfa}, while NNLO+PS predictions for the full VH process have been achieved in Refs.~\cite{Astill:2016hpa,Astill:2018ivh, Alioli:2019qzz}. Matched predictions including the $H\to b \bar{b}$ decay at full NNLO+PS accuracy have been obtained in Ref.~\cite{Zanoli:2021iyp}.

\subsection{Inclusive Cross Section Predictions}
\label{sec:vh_inclxs_breakdown}

The results compiled in Tables~\ref{tab:w+h_xs_7tev}-\ref{tab:zh_xs_14tev} are obtained from the following combination of the different contributions to the cross section
\begin{align}
    \sigma^{VH} &= \sigma_\mathrm{DY}^{VH}(1+ \delta_\mathrm{EW})+ \sigma_\text{top} + \delta_{\tiny VZ}~\sigma^\mathrm{NLO}_{ggZH},
    \label{eq:sigma_vh_breakdown}
\end{align}
where $\sigma_\mathrm{DY}^{VH}$ is the Drell-Yan-like contribution at NNLO QCD, $\delta_\mathrm{EW}$ is a multiplicative factor that accounts for the relative EW corrections at NLO, $\delta_{\tiny VZ}$ equals 1 when $V=Z$ and 0 otherwise, and $\sigma^\mathrm{NLO}_{ggZH}$ is the NLO QCD cross section for the $gg$-initiated channel. The term $\sigma_\text{top}$ is defined as
\begin{align}
     \sigma_\text{top} &= \sigma_\mathrm{I}^{VH} + \delta_{VZ}~\sigma_\mathrm{II}^{ZH},
\end{align}
where $\sigma_\mathrm{I}^{VH}$ and $\sigma_\mathrm{II}^{ZH}$ are contributions mediated by top-quark loops in the $q\bar{q}$-initiated channel.
The terms $\sigma^\mathrm{NLO}_{ggZH}$ and $\sigma_\mathrm{II}^{ZH}$  are present only in the $ZH$ final state. In the Tables~\ref{tab:w+h_xs_7tev}-\ref{tab:zh_xs_14tev} we provide results for on-shell bosons.

The results for all the contributions in Eq.~\eqref{eq:sigma_vh_breakdown} except $\sigma^\mathrm{NLO}_{ggZH}$ were always obtained using \texttt{vh@nnlo}~\cite{Harlander:2018yio}. The results for $\sigma^\mathrm{NLO}_{ggZH}$ at centre-of-mass (CM) energies $\sqrt{s}=\{13, 13.6, 14\}$~TeV are taken from Ref.~\cite{CampilloAveleira:2025rbh}, which includes NLO QCD corrections with top-quark mass effects at $\mathcal{O}(\%)$ accuracy. This reference provides a comparison of the virtual corrections reported in Refs.~\cite{Chen:2022rua,Degrassi:2022mro} and supplements the latter with the new results for the real corrections using \texttt{GoSam}~\cite{Braun:2025afl}. Results for values of the Higgs mass not present in Ref.~\cite{CampilloAveleira:2025rbh} have been obtained with a linear interpolation, which is accurate at $\mathcal{O}(10^{-4})$, given the mild dependence of the cross section on this parameter. For $\sqrt{s}=\{7, 8\}$~TeV, also the values of $\sigma^\mathrm{NLO}_{ggZH}$ are computed using \texttt{vh@nnlo}. 

The uncertainties due to the factorization and renormalization scale variation are denoted as $\Delta_\text{scale}$ and computed from a 7-point scale variation in the interval
\begin{align}
    M_{VH}/3 \leq \mu_F ,\mu_R \leq 3 M_{VH},
\end{align}
with $M_{VH}$ the invariant mass of the $VH$ system. Only in the case in which $\sigma^\mathrm{NLO}_{ggZH}$ is taken from Ref.~\cite{CampilloAveleira:2025rbh}, the scale uncertainties for the latter are obtained from a 7-point scale variation in the interval
\begin{align}
    M_{ZH}/4 \leq \mu_F ,\mu_R \leq M_{ZH}
\end{align}
and effects due to the different choice of scales are neglected. The PDF and $\alpha_s$ uncertainties (denoted as $\Delta_\text{PDF}$ ad $\Delta_{\alpha_s}$, respectively) are calculated following the recommendations in \url{https://twiki.cern.ch/twiki/bin/view/LHCPhysics/LHCHWG136TeVxsec}.



